%! Complex Numbers — Unit 2, Lesson 2.4 — Slides (Beamer)
\documentclass[aspectratio=169, 11pt]{beamer}
\usepackage{algebra2-beamer}   % provides \forestheader{} and \sectionlabel[color]{}
% Standard coordinate grid + axes: {xmin}{xmax}{ymin}{ymax}
\newcommand{\lgrid}[4]{%
  \draw[step=1, linegray!45, thin] (#1,#3) grid (#2,#4);
  \draw[->, thick, charcoal] (#1,0)--(#2,0) node[below right, font=\tiny]{$x$};
  \draw[->, thick, charcoal] (0,#3)--(0,#4) node[left, font=\tiny]{$y$};}

\begin{document}

% TITLE SLIDE (CourseName is not defined in beamer, so the name is literal)
{
\setbeamercolor{background canvas}{bg=forest}
\begin{frame}[plain]
  \vspace{2.8cm}\hspace{0.55cm}%
  \begin{minipage}{0.9\textwidth}
    {\color{goldacc}\bfseries\small LESSON 2.4}\\[0.5cm]
    {\color{white}\bfseries\LARGE Complex Numbers}\\[1.2cm]
    {\color{forestmid}\small Algebra 2: Shepherd~$\cdot$~Unit 2: Quadratic Functions}
  \end{minipage}
\end{frame}
}

% HOOK — the wall
\begin{frame}[t, plain]
  \forestheader{The wall: $x^2=-4$}
  \sectionlabel{Hook}
  \begin{columns}[T]
    \begin{column}{0.55\textwidth}
      In 2.3 we solved everything \emph{except} equations like
      \[ x^2=-4. \]
      \begin{itemize}\setlength{\itemsep}{4pt}
        \item No \emph{real} number squares to a negative.
        \item $y=x^2+4$ never crosses the $x$-axis --- no real zeros.
      \end{itemize}
      \textbf{Don't stop at ``no solution.'' Build a number that has one.}
    \end{column}
    \begin{column}{0.45\textwidth}
      \centering
      \begin{tikzpicture}[scale=0.5]
        \lgrid{-3}{3}{-1}{7}
        \begin{scope}\clip (-3,-1) rectangle (3,7);
          \draw[very thick, forest, domain=-1.7:1.7, samples=70, smooth] plot (\x,{\x*\x+4});
        \end{scope}
        \fill[charcoal] (0,4) circle (3pt) node[right, font=\tiny]{$(0,4)$};
      \end{tikzpicture}
    \end{column}
  \end{columns}
\end{frame}

% THE IMAGINARY UNIT
\begin{frame}[t, plain]
  \forestheader{The imaginary unit $i$}
  \sectionlabel[goldacc]{Definition}
  \begin{columns}[T]
    \begin{column}{0.52\textwidth}
      \begin{block}{Definition}
        $i=\sqrt{-1}, \qquad\text{so}\qquad i^2=-1.$
      \end{block}
      One rule, $i^2=-1$, powers everything.
      \[ \sqrt{-4}=\sqrt4\cdot\sqrt{-1}=2i. \]
    \end{column}
    \begin{column}{0.48\textwidth}
      \textbf{The wall opens:}
      \[ x^2=-4 \ \Rightarrow\ x=\pm\sqrt{-4}=\pm2i. \]
      \vspace{0.3cm}
      \begin{block}{Check}
        $(2i)^2=4i^2=4(-1)=-4.$ \checkmark
      \end{block}
    \end{column}
  \end{columns}
\end{frame}

% NEGATIVE RADICALS + a+bi
\begin{frame}[t, plain]
  \forestheader{Negative radicals \& standard form $a+bi$}
  \sectionlabel{Rewrite}
  \begin{columns}[T]
    \begin{column}{0.5\textwidth}
      For $k>0$: \ $\sqrt{-k}=i\sqrt{k}$ \ (pull $i$ out \emph{first}).
      \[ \sqrt{-9}=3i, \qquad \sqrt{-50}=5i\sqrt2. \]
      \begin{block}{Trap}
        $\sqrt{-2}\cdot\sqrt{-2}=2i^2=-2$, \ \textbf{not} $\sqrt4=2$.
      \end{block}
    \end{column}
    \begin{column}{0.5\textwidth}
      \textbf{Standard form} $a+bi$:
      \begin{itemize}\setlength{\itemsep}{3pt}
        \item $a=$ real part, $b=$ imaginary part.
        \item $b=0\Rightarrow$ real; $a=0\Rightarrow$ pure imaginary.
        \item $7-\sqrt{-16}=7-4i$ \ ($a=7,\ b=-4$).
      \end{itemize}
    \end{column}
  \end{columns}
\end{frame}

% OPERATIONS
\begin{frame}[t, plain]
  \forestheader{Add, subtract, multiply}
  \sectionlabel[goldacc]{Operate}
  \begin{columns}[T]
    \begin{column}{0.5\textwidth}
      \textbf{Add / subtract} --- combine like terms:
      \[ (3+2i)+(5-8i)=8-6i. \]
      \[ (4-i)-(6-3i)=-2+2i. \]
    \end{column}
    \begin{column}{0.5\textwidth}
      \textbf{Multiply} --- FOIL, then $i^2=-1$:
      \[ (2+3i)(4-i)=8+10i-3i^2=11+10i. \]
      \begin{block}{Conjugates are real}
        $(a+bi)(a-bi)=a^2+b^2$; \ $(5+2i)(5-2i)=29.$
      \end{block}
    \end{column}
  \end{columns}
\end{frame}

% POWERS OF i
\begin{frame}[t, plain]
  \forestheader{Powers of $i$ cycle}
  \sectionlabel{Pattern}
  \begin{columns}[T]
    \begin{column}{0.5\textwidth}
      \begin{block}{Period 4}
        \[ i^1=i,\ \ i^2=-1,\ \ i^3=-i,\ \ i^4=1,\ \ i^5=i,\dots \]
      \end{block}
      \textbf{Shortcut:} reduce the exponent \emph{mod} $4$.
    \end{column}
    \begin{column}{0.5\textwidth}
      \[ i^{23}:\ 23\bmod4=3 \ \Rightarrow\ i^3=-i. \]
      \[ i^{40}:\ 40\bmod4=0 \ \Rightarrow\ i^4=1. \]
      \footnotesize A multiple of $4$ always lands back on $1$.
    \end{column}
  \end{columns}
\end{frame}

% CAPSTONE + NEXT
\begin{frame}[t, plain]
  \forestheader{Where complex numbers come from \& what's next}
  \sectionlabel{Summary}
  \begin{columns}[T]
    \begin{column}{0.55\textwidth}
      \begin{block}{Complete the square on $x^2+2x+5=0$}
        \[ (x+1)^2=-4 \ \Rightarrow\ x=-1\pm2i. \]
      \end{block}
      The first quadratic with \emph{complex} solutions --- and they're conjugates.
    \end{column}
    \begin{column}{0.45\textwidth}
      \begin{block}{Next}
        \textbf{2.5 Quadratic Formula \& Discriminant}
        \[ x=\frac{-b\pm\sqrt{b^2-4ac}}{2a}. \]
        The sign of $b^2-4ac$ predicts real vs.\ complex roots.
      \end{block}
    \end{column}
  \end{columns}
\end{frame}

\end{document}
